\documentclass[11pt]{article}
\usepackage{amsmath}
\usepackage{amsfonts}
\usepackage{amssymb}
\usepackage{geometry}
\usepackage{graphicx}
\usepackage{algorithm}
\usepackage{algorithmic}

\geometry{margin=1in}

\title{Stability Analysis and Minimal Tau Estimation for Discrete-Time Dynamical Systems}
\author{Stability Analysis Documentation}
\date{\today}

\begin{document}

\maketitle

\section{Introduction}

This document explains the mathematical method for estimating the minimal time constant $\tau$ that guarantees stability of a discrete-time dynamical system. The system is characterized by a weight matrix $\mathbf{W}$ and transformed into a system matrix $\mathbf{A}$ through a specific parametrization involving $\tau$.

\section{System Definition}

Consider a discrete-time dynamical system defined by the system matrix:
\begin{equation}
\mathbf{A} = \mathbf{I} + \frac{\mathbf{W} - \mathbf{I}}{\tau}
\end{equation}

where:
\begin{itemize}
    \item $\mathbf{W}$ is the $n \times n$ weight matrix
    \item $\mathbf{I}$ is the $n \times n$ identity matrix
    \item $\tau > 0$ is the time constant parameter
    \item $\mathbf{A}$ is the resulting system matrix
\end{itemize}

\section{Stability Condition}

For a discrete-time linear system to be stable, all eigenvalues of the system matrix $\mathbf{A}$ must lie within the unit circle in the complex plane:

\begin{equation}
|\lambda_{\mathbf{A}}| < 1 \quad \text{for all eigenvalues } \lambda_{\mathbf{A}} \text{ of } \mathbf{A}
\end{equation}

This condition ensures that the system's response decays over time rather than growing unboundedly.

\section{Eigenvalue Relationship}

\subsection{Transformation of Eigenvalues}

If $\lambda_{\mathbf{W}}$ is an eigenvalue of $\mathbf{W}$ with corresponding eigenvector $\mathbf{v}$, then:
\begin{align}
\mathbf{W}\mathbf{v} &= \lambda_{\mathbf{W}}\mathbf{v}
\end{align}

The corresponding eigenvalue of $\mathbf{A}$ is:
\begin{align}
\mathbf{A}\mathbf{v} &= \left(\mathbf{I} + \frac{\mathbf{W} - \mathbf{I}}{\tau}\right)\mathbf{v}\\
&= \mathbf{v} + \frac{\mathbf{W}\mathbf{v} - \mathbf{v}}{\tau}\\
&= \mathbf{v} + \frac{\lambda_{\mathbf{W}}\mathbf{v} - \mathbf{v}}{\tau}\\
&= \left(1 + \frac{\lambda_{\mathbf{W}} - 1}{\tau}\right)\mathbf{v}
\end{align}

Therefore, the eigenvalue relationship is:
\begin{equation}
\lambda_{\mathbf{A}} = 1 + \frac{\lambda_{\mathbf{W}} - 1}{\tau}
\end{equation}

\section{Stability Constraint Derivation}

\subsection{Complex Eigenvalue Analysis}

Let $\lambda_{\mathbf{W}} = a + bi$ where $a, b \in \mathbb{R}$ are the real and imaginary parts respectively.

The corresponding eigenvalue of $\mathbf{A}$ becomes:
\begin{align}
\lambda_{\mathbf{A}} &= 1 + \frac{(a + bi) - 1}{\tau}\\
&= 1 + \frac{a - 1}{\tau} + \frac{bi}{\tau}\\
&= \left(1 + \frac{a - 1}{\tau}\right) + \frac{b}{\tau}i
\end{align}

\subsection{Magnitude Constraint}

For stability, we require $|\lambda_{\mathbf{A}}|^2 < 1$:
\begin{align}
|\lambda_{\mathbf{A}}|^2 &= \left(1 + \frac{a - 1}{\tau}\right)^2 + \left(\frac{b}{\tau}\right)^2 < 1
\end{align}

Expanding:
\begin{align}
\left(1 + \frac{a - 1}{\tau}\right)^2 + \frac{b^2}{\tau^2} &< 1\\
1 + 2\frac{a - 1}{\tau} + \frac{(a - 1)^2}{\tau^2} + \frac{b^2}{\tau^2} &< 1\\
2\frac{a - 1}{\tau} + \frac{(a - 1)^2 + b^2}{\tau^2} &< 0\\
2\frac{a - 1}{\tau} + \frac{|\lambda_{\mathbf{W}} - 1|^2}{\tau^2} &< 0
\end{align}

Since $|\lambda_{\mathbf{W}} - 1|^2 = (a - 1)^2 + b^2$.

\subsection{Solving for $\tau$}

Multiplying through by $\tau^2$ (assuming $\tau > 0$):
\begin{equation}
2(a - 1)\tau + |\lambda_{\mathbf{W}} - 1|^2 < 0
\end{equation}

Rearranging:
\begin{equation}
2(a - 1)\tau < -|\lambda_{\mathbf{W}} - 1|^2
\end{equation}

For this inequality to have a solution, we need $a - 1 < 0$, i.e., $a < 1$.

\textbf{Critical Observation}: If any eigenvalue of $\mathbf{W}$ has real part $\geq 1$, the system cannot be stabilized for any finite $\tau$.

When $a < 1$, we have $a - 1 < 0$, so:
\begin{align}
\tau &> \frac{-|\lambda_{\mathbf{W}} - 1|^2}{2(a - 1)}\\
\tau &> \frac{|\lambda_{\mathbf{W}} - 1|^2}{2(1 - a)}
\end{align}

\section{Algorithm for Minimal $\tau$ Estimation}

\begin{algorithm}
\caption{Minimal Tau Estimation}
\begin{algorithmic}[1]
\STATE \textbf{Input:} Weight matrix $\mathbf{W}$
\STATE \textbf{Output:} Minimal $\tau$ or indication if system cannot be stabilized

\STATE Compute eigenvalues $\{\lambda_1, \lambda_2, \ldots, \lambda_n\}$ of $\mathbf{W}$
\STATE Initialize $\tau_{constraints} = \{\}$

\FOR{each eigenvalue $\lambda_i = a_i + b_i i$}
    \IF{$a_i \geq 1$}
        \RETURN "System cannot be stabilized"
    \ENDIF
    
    \STATE Compute $\tau_{min,i} = \frac{|\lambda_i - 1|^2}{2(1 - a_i)}$
    \STATE Add $\tau_{min,i}$ to $\tau_{constraints}$
\ENDFOR

\STATE $\tau_{minimal} = \max(\tau_{constraints})$
\RETURN $\tau_{minimal}$
\end{algorithmic}
\end{algorithm}

\section{Implementation Details}

The constraint for each eigenvalue $\lambda_{\mathbf{W}} = a + bi$ is:
\begin{equation}
\tau > \frac{|\lambda_{\mathbf{W}} - 1|^2}{2(1 - a)}
\end{equation}

where:
\begin{itemize}
    \item $|\lambda_{\mathbf{W}} - 1|^2 = (a - 1)^2 + b^2$
    \item The constraint is valid only when $a < 1$
    \item The most restrictive constraint (largest $\tau$ requirement) determines the overall minimal $\tau$
\end{itemize}

\section{Physical Interpretation}

The parameter $\tau$ can be interpreted as a time constant that controls the rate of system evolution. Larger values of $\tau$ lead to:
\begin{itemize}
    \item Slower system dynamics
    \item More stable behavior
    \item Eigenvalues of $\mathbf{A}$ closer to 1
\end{itemize}

The minimal $\tau$ represents the smallest time constant that prevents the system from exhibiting exponential growth, ensuring all transients decay over time.

\section{Practical Considerations}

\subsection{Safety Margin}
In practice, it's recommended to use:
\begin{equation}
\tau_{practical} = \alpha \cdot \tau_{minimal}
\end{equation}
where $\alpha > 1$ (typically $\alpha = 1.1$ to $1.5$) provides a safety margin.

\subsection{Numerical Stability}
When implementing the algorithm:
\begin{itemize}
    \item Use appropriate numerical precision for eigenvalue computation
    \item Handle near-zero denominators carefully
    \item Consider eigenvalues with real parts very close to 1 as potentially problematic
\end{itemize}

\section{Conclusion}

The minimal $\tau$ estimation provides a rigorous mathematical framework for ensuring system stability. The method:
\begin{enumerate}
    \item Analyzes all eigenvalues of the weight matrix $\mathbf{W}$
    \item Derives individual stability constraints for each eigenvalue
    \item Selects the most restrictive constraint as the minimal $\tau$
    \item Guarantees that $\tau > \tau_{minimal}$ results in a stable system
\end{enumerate}

This approach is both theoretically sound and computationally efficient, making it suitable for practical stability analysis of dynamical systems.

\end{document} 